\documentclass[10pt,a4paper]{amsart}
\usepackage[margin=19mm]{geometry}
\usepackage{amsmath,amssymb,mathtools}
\usepackage[scr=rsfs,cal=euler]{mathalfa}
\usepackage{xfrac}
\usepackage[unicode,hidelinks,bookmarks=false]{hyperref}
\newcommand{\ie}{i.e.\ }
\newcommand{\cf}{cf.\ }
\newcommand{\resp}{respectively\ }
\newcommand{\Subsection}{\subsection}
\providecommand{\nobreakdash}{\nobreak\hbox{-}}
\setlength{\emergencystretch}{2em}
\multlinegap0pt
\usepackage{amssymb}
\newtheorem{lemma}{Lemma}
\newcommand{\DD}{\mathbb D}
\newcommand{\GG}{\mathbb G_2}
\title{Isometries in the symmetrized bidisc, II}
\author{Armen Edigarian}
\date{Source: arXiv:2608.08461v2 (CC BY 4.0)}
\begin{document}
\maketitle

\noindent\emph{Source and licence.} Isometries in the symmetrized bidisc, II. Version arXiv:2608.08461v2 (CC BY 4.0), distributed by arXiv under the Creative Commons Attribution 4.0 International licence, http://creativecommons.org/licenses/by/4.0/, as recorded in the arXiv licence metadata for this exact version and shown on the abstract page https://arxiv.org/abs/2608.08461v2. Full licence notice: \texttt{licenses/arxiv-2608.08461-CC-BY-4.0.txt}.\par\smallskip

\noindent\textit{Editorial scope.} Counted unit: the article's lemma lem:metric-derivative (source0.tex lines 150--155). The article's complete original proof of it is delivered with it (source0.tex lines 157--168, one \texttt{proof} environment of 12 lines). No mathematics is written, repaired or re-derived: the statement and the proof are the article's own lines, in the article's own notation, and no retained line is substituted or re-flowed.  No further source-local context is needed: every cross-reference of every retained line resolves inside the two delivered spans.  Nothing was omitted from any retained span, so the package carries no omission record.  No retained line contains a citation, so the wrapper carries no bibliography and no bibliography entry.  No retained line contains a figure environment, an image-inclusion command or drawing code, so no image is delivered and the package declares no asset dependency.  Editorial formatting only, in addition: an AMS \texttt{amsart} preamble that restates the article's own declarations and macros used by the retained lines, namely the environments \texttt{lemma} on separate counters, together with the standard packages those lines need.  The retained environments print in this document as Lemma 1 in order of appearance, whereas the article prints the same items with the numbering of its own full document.  The complete native source of the article as distributed in the arXiv e-print for this exact version is bundled unchanged as \texttt{sources/207162/source0.tex}.
\begin{lemma}\label{lem:metric-derivative}
If $q_t=q+tX+o(t)$ in $\GG$ as $t\to0$ through real values, then
\[
 \lim_{t\to0}\frac{k_{\GG}(q_t,q)}{|t|}=\kappa_{\GG}(q;X).
\]
\end{lemma}

\begin{proof}
For every $F\in\mathcal O(\GG,\DD)$, contractivity of the Carath\'eodory
distance gives
$\rho(F(q_t),F(q))\le k_{\GG}(q_t,q)$.  Dividing by $|t|$, passing to the
limit inferior and taking the supremum over $F$ gives
$\kappa_{\GG}(q;X)=\gamma_{\GG}(q;X)\le\liminf k_{\GG}(q_t,q)/|t|$.
For the reverse inequality join $q$ to $q_t$ by the straight segment.  For
small $t$ it lies in a compact subset of $\GG$, and continuity of
$\kappa_{\GG}$ gives an integrated length
$|t|\kappa_{\GG}(q;X)+o(|t|)$.  Since $k_{\GG}$ is the integrated distance of
$\kappa_{\GG}$, the opposite inequality follows.
\end{proof}

\end{document}
